\begin{definition}
\label{definition-quasi-projective}
\begin{reference}
\cite[II, Definition 5.3.1]{EGA} and \cite[page 103]{H}
\end{reference}
Let $f : X \to S$ be a morphism of schemes.
\begin{enumerate}
\item We say $f$ is {\it quasi-projective} if $f$ is of finite type
and there exists an $f$-relatively ample invertible $\mathcal{O}_X$-module.
\item We say $f$ is {\it H-quasi-projective} if there exists
a quasi-compact immersion $X \to \mathbf{P}^n_S$ over $S$ for some
$n$.\footnote{This is not exactly the same as the definition in Hartshorne.
Namely, the definition in Hartshorne (8th corrected printing, 1997) is that
$f$ should be the composition of an open immersion followed by a H-projective
morphism (see Definition \ref{definition-projective}), which does not imply
$f$ is quasi-compact. See
Lemma \ref{lemma-H-quasi-projective-open-H-projective} for
the implication in the other direction.}
\item We say $f$ is {\it locally quasi-projective} if there exists
an open covering $S = \bigcup V_j$ such that each $f^{-1}(V_j) \to V_j$
is quasi-projective.
\end{enumerate}
\end{definition}